% !TeX root = ../../Book4.tex
% 纲要标题：### 14.5 多重次调和与 Levi 几何
% 纲要位置：几何学计划纲要.md，第 4043 行。
% 正文状态：未撰写；本文件目前是框架。
% 本次扩充的范围、先修与边界：先修为本章多元全纯坐标、定义函数与 Levi 形式，以及 IV.13 的接触分布。固定复维数至少二的复欧氏空间中的光滑实超曲面，要求定义函数微分非零，构造复切 CR 分布并说明其形式可积性；完成球面、Heisenberg 超曲面及严格伪凸所诱导接触结构的计算。只处理嵌入超曲面模型，不承担抽象 CR 结构的一般可嵌入性、Kohn 理论或一般共形几何。

% 纲要内容开始
% 
% > 补充写作范围与外部输入：先修为本章多元全纯坐标、定义函数与 Levi 形式，以及 IV.13 的接触分布。固定复维数至少二的复欧氏空间中的光滑实超曲面，要求定义函数微分非零，构造复切 CR 分布并说明其形式可积性；完成球面、Heisenberg 超曲面及严格伪凸所诱导接触结构的计算。只处理嵌入超曲面模型，不承担抽象 CR 结构的一般可嵌入性、Kohn 理论或一般共形几何。
% 
% * Levi 形式及 psh 函数
% * 严格 psh 与局部势
% * 伪凸性和域的整体问题
% * 与后续带权 ∂̄ 估计的连接
% 
% * 实超曲面的 CR 分布与形式可积性
% * 球面与 Heisenberg 超曲面的 Levi 形式
% * 严格伪凸超曲面的接触结构
% 
% 
% ---
% 
% 纲要内容结束

\section{多重次调和与 Levi 几何}
\label{g4:sec:1405}

待撰写

\subsection{Levi 形式及 psh 函数}
\label{g4:subsec:1405:01}
待撰写

\subsection{严格 psh 与局部势}
\label{g4:subsec:1405:02}
待撰写

\subsection{伪凸性和域的整体问题}
\label{g4:subsec:1405:03}
待撰写

\subsection{\texorpdfstring{与后续带权 \(\bar{\partial}\) 估计的连接}{与后续带权 \={∂} 估计的连接}}
\label{g4:subsec:1405:04}
待撰写

\subsection{实超曲面的 CR 分布与形式可积性}
\label{g4:subsec:1405:05}
待撰写

\subsection{球面与 Heisenberg 超曲面的 Levi 形式}
\label{g4:subsec:1405:06}
待撰写

\subsection{严格伪凸超曲面的接触结构}
\label{g4:subsec:1405:07}
待撰写

% 纲要条目排为编号小节标题，并进入目录与书签；说明段仍保留于纲要注释中。
